%%%%%%%%%%%%%%%%%%%%%%%%%%%%%%%%%%%%%%%%%
% Friggeri Resume/CV
% XeLaTeX Template
% Version 1.0 (5/5/13)
%
% This template has been downloaded from:
% http://www.LaTeXTemplates.com
%
% Original author:
% Adrien Friggeri (adrien@friggeri.net)
% https://github.com/afriggeri/CV
%
% License:
% CC BY-NC-SA 3.0 (http://creativecommons.org/licenses/by-nc-sa/3.0/)
%
% Important notes:
% This template needs to be compiled with XeLaTeX and the bibliography, if used,
% needs to be compiled with biber rather than bibtex.
%
%%%%%%%%%%%%%%%%%%%%%%%%%%%%%%%%%%%%%%%%%

\documentclass[print]{friggeri-cv} % Add 'print' as an option into the square bracket to remove colors from this template for printing

\addbibresource{bibliography.bib} % Specify the bibliography file to include publications

\begin{document}

\header{Johan Andreas}{ Versteegh}{Problem Solver, Communicator, Craftsman} % Your name and current job title/field
%----------------------------------------------------------------------------------------
%	SIDEBAR SECTION
%----------------------------------------------------------------------------------------

\begin{aside} % In the aside, each new line forces a line break
\section{Birthdate}
10-october-1972
\section{contact}
Koppelsedijk 18
4191 LC
The Nederlands
%Nederland
~
Mob: +31(0)627042986
~
\href{mailto:johan@techplaygroup.com}{johan@techplaygroup.com}
\href{http://www.techplaygroup.com}{www.techplaygroup.com}
\section{languages}
Dutch
English
Norwegian
German
French
\section{skills}
{Sensor Technology
Mechatronics
System Control:
(Labview, Beckhoff)
Embedded Systems:
ASM, C
3D design}
\section{interests}
emerging technologies
anomaly detection
teaching
\end{aside}

%----------------------------------------------------------------------------------------
%	Summary
%----------------------------------------------------------------------------------------

%My professional life passed through several paths:
%(1) From an early age is was interested in technology and absorbing information came easy to me %(2) I'm attracted to aesthetics and art, and the highlight was my Film&TV studies, the short films I made and the awards I got for my film Tukim;
%(3) I found out in my twenties %that I can use my creativity for Business Development, and during my years in ATS, I reached outstanding sales achievements with the Phantom product line in Israel.

%----------------------------------------------------------------------------------------
%	EDUCATION SECTION
%----------------------------------------------------------------------------------------

\section{education}

\begin{entrylist}
%------------------------------------------------
\entry
{2001--2002}
%{Elektrotechniek {\normalfont (parttime)}}
{Engineering and Life Sciences {\normalfont (parttime)}}
% {Hogeschool van Arnhem en Nijmegen}
{HAN University of Applied Sciences, Arnhem}
{Automotive Engineering}
%------------------------------------------------
\entry
{1993--1998}
%{Industrieel Ontwerpen {\normalfont (propedeuse)}}
{Industrial Design {\normalfont (No Degree)}}
{TU Delft}
{Material sciences, ergonomics, mechanical and electrical engineering, Human machine interaction}
%{Industrial design is the use of both applied art and applied science to improve the aesthetics, ergonomics, functionality, and/or usability of a product, and it may also be used to improve the product's marketability and even production. The role of an industrial designer is to create and execute design solutions for problems of form, usability, physical ergonomics, marketing, brand development, and sales.}
%------------------------------------------------
\entry
%1992- 1994    Mijnbouwkunde, TUDelft.
{1992--1994}
{Applied Earth Sciences {\normalfont (propedeuse)}}
{TU Delft}
{Petroleum Engineering and Geosciences, Geoscience and Remote Sensing}
\end{entrylist}

%----------------------------------------------------------------------------------------
%	WORK EXPERIENCE SECTION
%----------------------------------------------------------------------------------------

\section{experience}

\begin{entrylist}
%------------------------------------------------
\entry
{2011--2013}
{INCAS³}
{Assen, The Netherlands}
{\emph{Various Functions:\\}
\lang{nl}{INCAS³ is een onafhankelijk non-profit research instituut dat zich bezig houdt met het oplossen van uitdagende problemen binnen de industrie en op sociaal technisch gebied.}
\lang{en}{\href{http://www.incas3.eu/}{INCAS³} is an independent, private, non-profit research institute dedicated to solving challenging industrial and social technological problems.}
\lang{nl}
{Functie:
Het ontwikkelen en inkopen van test opstellingen om de door de wetenschappers gevonden theorieën te bewijzen.
Interim Head of engineering.
Industrial Researcher: de schakel vormen tussen de wetenschap en de industrie en in dat proces de daaruit voortkomende (deel)oplossingen proberen om te zetten naar produceerbare eindproducten.
Projecten:
SPRINT, SMARTBOT, SENSOR CITY, MoreWise. 
Binnen deze projecten was een van mijn taken de Phd’s en Master studenten op technisch vlak te begeleiden.}
\lang{en}{
\\Functions and Responsibilities:
\begin{itemize}
	\item Industrial Researcher
	\begin{itemize}
		\item Designed and developed test setups to prove scientific theories.
		\item Forming links between science and industry and in this process transform (part of) the solutions into products ready for manufacturing.
		\lang{nl}{\item SPRINT, SMARTBOT, SENSOR CITY, MoreWise. Binnen deze projecten was een van mijn taken de Phd’s en Master studenten op technisch vlak te begeleiden.}
		\lang{en}{\item Managing and Engineering in Dutch and European collaborations:
		\begin{itemize}
			\item \href{http://www.hannn.eu/nl/nieuws/archief/center-of-research-excellence-sprint}{SPRINT}
			\item \href{http://www.smartbot.eu/}{SMARTBOT}
			\item \href{http://www.incas3.eu/focus-areas/projects/morewise}{MoreWise}
			\end{itemize}
		\item Supported Phd’s and Master students involved in the projects.}
	\end{itemize}
	\item Interim Head of engineering
	\begin{itemize}
		\item Getting the engineering department up to date on the future in production and sensor technology.
%		\item Sourcing equipment and materials for new projects.
	\end{itemize}
		\item CTO, Otigy (Daughter company of INCAS³)
	\begin{itemize}
		\item Professionalized a functional prototype developed at INCAS³ and made it ready for mass production. 
		\item Managed production of the first 400 prototypes installed in a citywide sensor network: \href{http://www.sensorcity.nl}{(www.sensorcity.nl)\\}.
	\end{itemize}
\end{itemize}
}}
\end{entrylist}
%------------------------------------------------
\begin{entrylist}
\entry
{2006--2011}
{Technologies88}
{Amsterdam, The Netherlands}
{\emph{CTO\\}
\lang{nl}{Technologies88 specialiseerde zich in het ontwikkelen van innovatie gereedschappen die hulp bieden bij het nemen van kritieke beslissingen binnen sport, gezondheid en industrie.
Functies: 
Ontwikkelen elektronica en low level software.
Keuze en inkoop van passende componenten.
Ontwikkelen van test systemen voor ontwikkelde hardware.
Eruit voortvloeiende patenten uitwerken en indienen.
Behuizing ontwerpen en productie begeleiden.
Begeleiden massa productie.}
\lang{en}{Technologies88 specialized in the development of innovative technology to aid processes in the fields of Sport, Health and Industry where critical decisions need to be made in realtime.\\In 2010 a fully embedded golf training device was launched: \href{https://www.youtube.com/watch?v=774XEkbOrPo}{The Digital Golf Coach}. \href{https://www.youtube.com/watch?v=774XEkbOrPo}{The Digital Golf Coach} used an 3axis accelerometer to analyze the Golf swing and compare it to your personal best.}
\lang{en}{
\\Functions and Responsibilities:
\begin{itemize}
	\item Develop Electronics and Low level software.
	\item Sourcing components and materials.
	\item Develop test system and protocols for developed hardware.
	\item Write and manage patent applications for found innovations.
	\item Make product ready for production (Design enclosure, User Manual, Production tools).
	\item Manage production in Europe and China.\\
\end{itemize}
}}
\end{entrylist}
%------------------------------------------------
\begin{entrylist}
\entry
{2003--2005}
{Wireless Playground b.v.}
{Rotterdam, The Netherlands}
{\emph{Co-Founder} \\
\lang{nl}{Medeoprichter van een netwerk organisatie voor ontwikkeling en introductie van nieuwe producten.
Functie:  Ontwikkeling van proof of concept en coördinatie van massa productie.
Partners: Skept, Philips RF, Elan Microchip Co., Intersolve.
}
\lang{en}{Wireless Playground was founded as a vehicle to develop and introduce new gadgets.\\
In 2004 the WiPla RDS module was introduced, a battery powered sound module remotely controllable via RDS.\\
This module enabled wireless control throughout most of Europe via cheap FM broadcasts.
}
\lang{en}{
\\Functions and Responsibilities:
\begin{itemize}
	\item Designing proof of concept and making it ready for mass production.
	\item Partners: Philips RF, Elan Microchip Co, Sanyo.\\
\end{itemize}
}
}

\entry
{2002--2014}
{Technology Playgroup Europe B.V.}
{Dieren, The Netherlands}
{\emph{Founder} \\
\lang{nl}{Dochter onderneming van Technology Playgroup Inc, opgericht om zich beter te kunnen richten op de Europese markt, en een betere infrastructuur te vormen voor technische ondersteuning.
Samen met TPG inc. zijn er vele producten ontwikkeld voor de film industrie
Patents: http://ip.com/patent/US7862522
}
\lang{en}{European daughter company of Technology Playgroup Inc, founded to better serve the European market.\\
Together with Technology Playgroup Inc. many products where developed for the Film and the Medical industry. 
\begin{itemize}
	\item \href{http://davebarclay.com/outabody.html}{Sensorbox}
	\item \href{http://davebarclay.com/outabody.html} {DataGlove}
	\item \href{http://conceptoverdrive.com/products/cancoder_.php} {CanCoder}\\
\end{itemize}
}
}

\entry
{2001--2002}
{Technology Playgroup Inc.}
{Montreal, Canada}
{\emph{Team leader Development team} \\
\lang{nl}{Ontwikkeling en productie van een draadloos sensor systeem voor interactieve media en bewegingsanalyse.
Verantwoordelijk voor de integratie van nieuwe technologische ontwikkelingen. 
}
\lang{en}{Technology Playgroup Inc. developed of a wireless sensor system for motion capture and interactive media using the latest sensor technology.\\ 
}
}

\lang{th}{
\entry
{1997--2000}
{More Stage Services}
{The Hague, The Netherlands}
{\emph{Theatre Engineer} \\
\lang{nl}{Geluid en licht technicus voor theater shows en evenementen.
Verantwoordelijk voor reparaties, onderhoud aan apparatuur en het ontwerpen en inkopen van multimedia installaties. 
}
\lang{en}{Sound and Light Designer/engineer, Repairs, maintenance and acquisition of multimedia installations.\\
}
}

\entry
{1998--2000}
{Haine Audio}
{Rotterdam, The Netherlands}
{\emph{Media System Engineer} \\
\lang{nl}{Multimedia/Netwerk technicus.
Verantwoordelijk voor het ontwerpen en analyseren van video, geluid, licht en computer netwerken op grote evenementen en het inkopen van nieuwe apparatuur. 
}
\lang{en}{Design and Analysis of multimedia systems on large events and acquisition of new technology.
}
}
}
%------------------------------------------------
\end{entrylist}

%----------------------------------------------------------------------------------------
%	PRojects
%----------------------------------------------------------------------------------------

\lang{th}{
\section{projects}

\begin{entrylist}
%------------------------------------------------
\entry
{2001--2002}
{Krisztina de Châtel \& Geert Mul}
{Rotterdam, The Netherlands}
{
\lang{nl}{Ontwikkeling en productie van beweging sensoren en integratie van 3D projecties voor de theater productie DYNAMIX.}
\lang{en}{Developing and integrating a motion tracking system for a moving stage with 3D projections for the theatre piece \href{http://www.kdechatel.com/english/krisztina-de-chatel/choreographies/dynamix/}{DYNAMIX}.}\\
}
%------------------------------------------------
\entry
{1999--Dec}
{BOANER \& De Schouwburg Rotterdam}
{Rotterdam The Netherlands}
{
\lang{nl}{Technische Productie van het X-samples Festival.}
\lang{en}{Technical Production of the X-samples Festival.}
}
%------------------------------------------------
\entry
{1999--Okt}
{CROSSFAIR – DANCE AND TECHNOLOGY SYMPOSIUM}
{Essen, Germany}
{
\lang{nl}{Technische Productie samen met SHINKANSEN, Londen.}
\lang{en}{Technical Production with SHINKANSEN,London.}
}
%------------------------------------------------
\entry
{1999--Mai}
{koerper.technik//body.technology}
{Berlin, Germany}
{
\lang{nl}{Artiest}
\lang{en}{Visiting artist and Technical Production.}
}

%------------------------------------------------
\end{entrylist}
}
%----------------------------------------------------------------------------------------
%	Patents SECTION
%----------------------------------------------------------------------------------------

\section{patents}

\begin{itemize}
	\item \href{http://ip.com/pat/WO2009034189A1}{WO2009034189A1, Training apparatus}
	\item \href{http://ip.com/pat/WO2011139153A1}{WO2011139153A1, Device and method for motion capture and analysis}
	\item \href{http://ip.com/patent/US7862522}{US7862522, Sensor glove}
\end{itemize}



\lang{nl}{
%----------------------------------------------------------------------------------------
%	COMMUNICATION SKILLS SECTION
%----------------------------------------------------------------------------------------

\section{communication skills}

\begin{entrylist}
%------------------------------------------------
\entry
{2011}
{Oral Presentation}
{California Business Conference}
{Presented the research I conducted for my Masters of Commerce degree.}
%------------------------------------------------
\entry
{2010}
{Poster}
{Annual Business Conference, Oregon}
{As part of the course work for BUS320, I created a poster analyzing several local businesses and presented this at a conference.}
%------------------------------------------------
\end{entrylist}


%----------------------------------------------------------------------------------------
%	INTERESTS SECTION
%----------------------------------------------------------------------------------------

\section{interests}

\textbf{professional:} data analysis, company profiling, risk analysis, economics, web design, web app creation, software design, marketing \textbf{personal:} piano, chess, cooking, dancing, running

%----------------------------------------------------------------------------------------
%	PUBLICATIONS SECTION
%----------------------------------------------------------------------------------------

\section{publications}

\printbibsection{article}{article in peer-reviewed journal} % Print all articles from the bibliography

\printbibsection{book}{books} % Print all books from the bibliography

\begin{refsection} % This is a custom heading for those references marked as "inproceedings" but not containing "keyword=france"
\nocite{*}
\printbibliography[sorting=chronological, type=inproceedings, title={international peer-reviewed conferences/proceedings}, notkeyword={france}, heading=subbibliography]
\end{refsection}

\begin{refsection} % This is a custom heading for those references marked as "inproceedings" and containing "keyword=france"
\nocite{*}
\printbibliography[sorting=chronological, type=inproceedings, title={local peer-reviewed conferences/proceedings}, keyword={france}, heading=subbibliography]
\end{refsection}

\printbibsection{misc}{other publications} % Print all miscellaneous entries from the bibliography

\printbibsection{report}{research reports} % Print all research reports from the bibliography

%----------------------------------------------------------------------------------------

%Multilanguage test
% Sections, subsections, subsubsections:

\sectionlang{no}{norksidoo} % in the first parameter you can freely use
                            % both short or long language codes
\sectionlang{en}{Section 1}
\sectionlang{nl}{Sectie 1}

% Strings that are displayed only in a certain language
\lang{nl}{hier dus nederlands.}
\lang{en}{This string shows up only in English.}
\lang{no}{Og denne på norsk.}
}


\end{document}